\lecture{22}{May 14}{}
Let \(X_1, X_2\) be continuous random variables in a probability space. Suppose we define 
\[
    Y_1 = g_1(X_1, X_2), \quad Y_2 = g_2(X_1, X_2),
\] 
where \(g_1, g_2: \mathbb{R} ^2 \to \mathbb{R}\). If for every \((y_1, y_2) \in \mathbb{R} ^2\), there is a unique \((x_1, x_2) \in \mathbb{R} ^2\) where \(x_1 = x_1(y_1, y_2), x_2 = x_2(y_1, y_2)\) s.t.
\[
    y_1 = g_1(x_1, x_2) \text{ and } y_2 = g_2(x_1, x_2), 
\]   
and 
\[
    J(x_1, x_2) = \left\vert \begin{pmatrix}
        \frac{\partial g_1}{\partial x_1}  & \frac{\partial g_1}{\partial x_2}   \\
         \frac{\partial g_2}{\partial x_1} & \frac{\partial g_2}{\partial x_2}   \\
    \end{pmatrix} \right\vert \neq 0,
\]
then 
\[
    f_{Y_1, Y_2}(y_1, y_2) = f_{X_1, X_2} \left( x_1 (y_1, y_2), x_2(y_1, y_2)\right) \left\vert J(x_1, x_2) \right\vert ^{-1}. 
\]
\begin{proof}
    Compute the joint cumulative distribution function for \(Y_1, Y_2\),
    \begin{align*}
        F_{Y_1, Y_2} (y_1, y_2) &= \mathbb{P} \left( Y_1 \le y_1, Y_2 \le y_2 \right) = \mathbb{P} \left( \left\{ (x_1, x_2): g_1(x_1, x_2) \le y_1, g_2(x_1, x_2) \le y_2 \right\}  \right) \\
        &= \iint_{\left\{ \substack{(x_1, x_2): g_1 \le y_1, g_2 \le y_2} \right\} } f_{X_1, X_2} (x_1, x_2) \, \mathrm{d} x_1 \, \mathrm{d} x_2.   
    \end{align*} 
    Now let \(u_1 = g_1(x_1, x_2), u_2 = g_2(x_1, x_2)\), we have 
    \begin{align*}
       \iint_{\left\{ \substack{(x_1, x_2): g_1 \le y_1, g_2 \le y_2} \right\} } f_{X_1, X_2} (x_1, x_2) \, \mathrm{d} x_1, \, \mathrm{d} x_2 &= \iint_{\left\{ u_1 \le y_1, u_2 \le y_2 \right\} } f_{X_1, X_2} \left( x_1(u_1, u_2), x_2(u_1, u_2) \right) \left\vert J(x_1, x_2) \right\vert^{-1} \, \mathrm{d} u_1 \, \mathrm{d} u_2.     
    \end{align*} 
    The density function is the partital distribition of \(y_1, y_2\). 
    \[
        f_{Y_1, Y_2} (y_1, y_2) = \frac{\partial^2 F_{Y_1, Y_2}}{\partial y_1 \partial y_2} = f_{X_1, X_2} (x_1, x_2) \vert J(x_1, x_2) \vert^{-1},  
    \] 
    where \(y_1 = g_1(x_1, x_2)\) and \(y_2 = g_2(x_1, x_2)\).  
\end{proof}


\begin{eg}
    \(Y_1 = X_1 + X_2, Y_2 = X_1 - X_2\), then 
    \[
        X_1 = \frac{1}{2} \left( Y_1 + Y_2 \right), \quad X_2 = \frac{1}{2} \left( Y_1 - Y_2 \right),  
    \] 
    and 
    \[
        J(x_1, x_2) = \left\vert \begin{pmatrix}
            1 & 1  \\
            1 & -1  \\
        \end{pmatrix} \right\vert = -2 \neq 0, 
    \]
    so 
    \[
        f_{Y_1, Y_2} (y_1, y_2) = f_{X_1, X_2} \left( \frac{y_1 + y_2}{2}, \frac{y_1 - y_2}{2} \right) \cdot \frac{1}{2}. 
    \]
\end{eg}

\begin{eg}
    Let \(X_1, X_2 \sim N(0, 1)\), independent. Let \(Y_1 = X_1 + X_2\) and \(Y_2 = X_1 - X_2\). What can be said about the joint distribution of \(Y_1, Y_2\)?    
\end{eg}
\begin{explanation}
    Recall that if \(X_1 \sim N(\mu _1, \sigma _1^2)\) and \(X_2 \sim N(\mu _2, \sigma _2^2)\), independent, then \(X_1 + X_2 \sim N(\mu _1 + \mu _2, \sigma _1^2 + \sigma _2^2)\), so \(Y_1 \sim N(0, 2)\) and \(Y_2 \sim N(0, 2)\). Thus, 
    \begin{align*}
        f_{X_1, X_2} (x_1, x_2) &= f_{X_1}(x_1) f_{X_2}(x_2) = \left( \frac{1}{\sqrt{2\pi } } e^{- \frac{x_2^2}{2}} \right) \left( \frac{1}{\sqrt{2 \pi} } e^{- \frac{x_2^2}{2}} \right) = \frac{1}{2 \pi } e^{- \frac{x_1^2 + x_2^2}{2}}.  
    \end{align*}     
    Thus, 
    \begin{align*}
        f_{Y_1, Y_2} (y_1, y_2) &= \frac{1}{2} f_{X_1, X_2} \left( \frac{y_1 + y_2}{2}, \frac{y_1 - y_2}{2} \right) = \frac{1}{4 \pi} e^{-\frac{1}{2} \left[ \left( \frac{y_1 + y_2}{2} \right)^2 + \left( \frac{y_1 - y_2}{2} \right)^2   \right] } \\
        &= \frac{1}{4 \pi} e^{-\frac{1}{2} \left[ \frac{y_1^2}{2} + \frac{y_2^2}{2} \right] } = \left( \frac{1}{\sqrt{2\pi \cdot 2} } e^{- \frac{y_1^2}{2 \cdot 2}} \right) \left( \frac{1}{\sqrt{2\pi \cdot 2} e^{-\frac{y_2^2}{2 \cdot 2}}} \right) = f_{Y_1}(y_1) f_{Y_2}(y_2),   
    \end{align*}
    so \(Y_1, Y_2\) are independent \(N(0, 2)\).  
\end{explanation}

\textbf{Fact:} If \(X_1, X_2\) are independent and identically distributed continuous random variables, and \(Y_1 = X_1 + X_2\) and \(Y_2 = X_1 - X_2\) are also independent, then \(X_1, X_2\) are normally distributed. 

More generally, given \(X_1, X_2, \dots , X_n\), if we define \(Y_i = g_i \left( X_1, X_2, \dots , X_n \right) \) s.t. 
\[
    (x_1, x_2, \dots , x_n) \mapsto (y_1, y_2, \dots , y_n)
\]  
is a bijection and 
\[
    J(x_1, \dots , x_n) = \det \left( \left( \frac{\partial g_i}{\partial x_j}  \right)_{i, j = 1}^n  \right) \neq 0, 
\]
then 
\[
    f_{Y_1, Y_2, \dots , Y_n} (y_1, \dots , y_n) = f_{X_1, X_2, \dots , X_n} \left( x_1, x_2, \dots , x_n \right) \left\vert J(x_1, x_2, \dots , x_n) \right\vert^{-1}.  
\]

\section{Expectation of functions of random variables}
Recall that if \(X\) is a discrete random variable taking non-negative integer values, then 
\[
    \mathbb{E} [X] = \sum_{n=1}^{\infty} \mathbb{P} (X \ge n), 
\] 
which is proved in the homework. 

\begin{lemma}
    Let \(X\) be a continuous random variable. 
    \[
        \mathbb{E} [X] = \int _0^{\infty} \mathbb{P} (X \ge t) \, \mathrm{d} t - \int _0^{\infty} \mathbb{P} (X \le -t) \, \mathrm{d} t. 
    \] 
\end{lemma}
\begin{proof}
    \begin{align*}
        \int _0^{\infty} \mathbb{P} (X \ge t) \, \mathrm{d} t &= \int_0^{\infty} \int _t^{\infty} f_X(x) \, \mathrm{d} x \, \mathrm{d} t = \int _0^{\infty} \int _0^x f_X(x) \, \mathrm{d} t \, \mathrm{d} x \\
        &= \int _0^{\infty} f_X(x) \int _0^x 1 \, \mathrm{d} t \, \mathrm{d} x = \int _0^{\infty} x f_X(x) \, \mathrm{d} x.         
    \end{align*}

    \begin{align*}
        \int _0^{\infty} \mathbb{P} (X \le -t) \, \mathrm{d} t &= \int _0^{\infty} \int _{-\infty }^{-t} f_X(x) \, \mathrm{d} x \, \mathrm{d} t = \int _{-\infty }^0 \int _0^{-x} f_X(x) \, \mathrm{d} t \, \mathrm{d} x \\
        &= -\int _{-\infty }^0 x f_X(x) \, \mathrm{d} x.      
    \end{align*}

    Hence,
    \[
        \int _0^{\infty} \mathbb{P} (X \ge t) \, \mathrm{d} t - \int _0^{\infty} \mathbb{P} (X \le -t) \, \mathrm{d} t = \int_{-\infty }^{\infty} x f_X(x) \, \mathrm{d} x = \mathbb{E} [X].    
    \]
\end{proof}

\begin{proposition}
    Let \(X_1, X_2\) be random variables on a probability space, and let \(g: \mathbb{R}^2 \to \mathbb{R} \), then 
    \begin{itemize}
        \item (discrete) \(\mathbb{E} [g(X_1, X_2)] = \sum_{x_1, x_2} g(x_1, x_2) p_{X_1, X_2} (x_1, x_2)\). 
        \item (continuous) \(\mathbb{E} [g(X_1, X_2)] = \iint_{\mathbb{R} ^2} g(x_1, x_2) f_{X_1, X_2}(x_1, x_2) \, \mathrm{d} x_1 \, \mathrm{d} x_2   \).
    \end{itemize}  
\end{proposition}

\begin{proof}
    Here we only prove the continuous setting. First, we have
    \begin{align*}
        \mathbb{E} [g(X_1, X_2)] &= \int _0^{\infty} \mathbb{P} (g(X_1, X_2) \ge t) \, \mathrm{d} t - \int _0^{\infty} \mathbb{P} (g(X_1, X_2) \le -t) \, \mathrm{d} t. 
    \end{align*}
    Now since 
    \begin{align*}
        \mathbb{P} (g(X_1, X_2) \ge t) = \iint_{\left\{ (x_1, x_2): g(x_1, x_2) \ge t \right\} } f_{X_1, X_2}(x_1, x_2) \, \mathrm{d} x_1 \, \mathrm{d} x_2,  
    \end{align*}
    we have 
    \begin{align*}
        \int _0^{\infty} \mathbb{P} (g(X_1, X_2) \ge t) \, \mathrm{d} t &= \int _0^{\infty} \iint_{\left\{ (x_1, x_2): g \ge t \right\} } f_{X_1, X_2} (x_1, x_2) \, \mathrm{d} x_1 \, \mathrm{d} x_2 \, \mathrm{d} t \\
        &= \iint_{\left\{ (x_1, x_2): g \ge 0 \right\} } \int _0^{g(x_1, x_2)} f_{X_1, X_2} (x_1, x_2) \, \mathrm{d} t \, \mathrm{d} x_1 \, \mathrm{d} x_2 \\
        &= \iint_{\left\{ (x_1, x_2): g \ge 0 \right\} } g(x_1, x_2) f_{X_1, X_2}(x_1, x_2) \, \mathrm{d} x_1 \, \mathrm{d} x_2.      
    \end{align*}
    Similarly,
    \begin{align*}
        \int_0^{\infty} \mathbb{P} (g(X_1, X_2) \le -t) \, \mathrm{d} t &= \int _0^{\infty} \iint_{\left\{ (x_1, x_2): g \le -t \right\} } f_{X_1, X_2} (x_1, x_2) \, \mathrm{d} x_1 \, \mathrm{d} x_2 \, \mathrm{d} t \\
        &= \iint_{\left\{ (x_1, x_2): g \le 0 \right\} } \int _0^{-g(x_1, x_2)} f_{X_1, X_2} (x_1, x_2) \, \mathrm{d} t \, \mathrm{d} x_1 \, \mathrm{d} x_2 \\
        &= - \iint_{\left\{ (x_1, x_2): g \le 0 \right\} } g(x_1, x_2) f_{X_1, X_2}(x_1, x_2) \, \mathrm{d} x_1 \, \mathrm{d} x_2.         
    \end{align*}
    Hence, 
    \begin{align*}
        \mathbb{E} [g(X_1, X_2)] &= \int_0^{\infty} \mathbb{P} (g(X_1, X_2) \ge t) \, \mathrm{d} t - \int_0^{\infty} \mathbb{P} (g(X_1, X_2) \le t) \, \mathrm{d} t \\
        &= \iint_{\left\{ (x_1, x_2): g \ge 0 \right\} } g(x_1, x_2) f_{X_1, X_2} (x_1, x_2) \, \mathrm{d} x_1 \mathrm{d} x_2 - \left( - \iint_{\left\{ (x_1, x_2): g \le 0 \right\} } g(x_1, x_2) f_{X_1, X_2}(x_1, x_2) \, \mathrm{d} x_1 \,\mathrm{d} x_2   \right) \\
        &= \iint_{\mathbb{R} ^2} g(x_1, x_2) f_{X_1, X_2} (x_1, x_2) \, \mathrm{d} x_1 \, \mathrm{d} x_2.       
    \end{align*}

\end{proof}

\begin{eg}
    A road of length \(L\). An accident occurs at portion \(X\), uniformly along the road. An ambulance is at portion \(Y\), uniformly along the road, independent of \(X\). How far must the ambulance travel to get to the accident? What is the expected distance?    
\end{eg}

\begin{explanation}
    We know \(X \sim \mathrm{Unif}([0, L]), Y \sim \mathrm{Unif}([0, L])  \) are independent, so 
    \[
        f_{X, Y}(x, y) = \begin{dcases}
            \frac{1}{L^2}, &\text{ if } 0 \le x, y \le L ;\\
            0, &\text{ otherwise} .
        \end{dcases}
    \]
    Now define \(\mathrm{dis}(x, y) = \vert x - y \vert  \), then 
    \begin{align*}
        \mathbb{E} [\mathrm{dis}(X, Y)] &= \iint_{\mathbb{R} ^2} \mathrm{dis}(x, y) f_{X, Y}(x, y) \, \mathrm{d} x \, \mathrm{d} y = \frac{1}{L^2} \int _0^L \int _0^L \vert x - y \vert \, \mathrm{d} x \, \mathrm{d} y \\
        &= \frac{1}{L^2} \int _0^L \left( \int _0^y (y - x) \, \mathrm{d} x + \int _y^L (x - y) \, \mathrm{d} x   \right) \, \mathrm{d} y \\
        &= \frac{1}{L^2} \int _0^L \left( \left[ xy - \frac{1}{2}x^2 \right]_0^y + \left[ \frac{1}{2}x^2 - xy \right]_y^L   \right) \, \mathrm{d} y \\
        &= \frac{1}{L^2} \int _0^L \frac{1}{2} y^2 + \frac{1}{2} L^2 - Ly + \frac{1}{2} y^2 \, \mathrm{d} y \\
        &= \frac{1}{L^2} \int _0^L y^2 + \frac{1}{2} L^2 - Ly \, \mathrm{d} y \\
        &= \frac{1}{L^2} \left[ \frac{1}{3} y^3 + \frac{1}{2} L^2 y - \frac{1}{2} L y^2 \right]_0^L \\
        &= \frac{1}{L^2} \left[ \frac{1}{3}L^3 + \frac{1}{2} L^3 - \frac{1}{2} L^3 \right] = \frac{L}{3}.             
    \end{align*}

\end{explanation}

\paragraph{Application (Linearity)} If \(g(X, Y) = aX + bY\), then 
\begin{align*}
    \mathbb{E} [aX + bY] &= \iint_{\mathbb{R} ^2} g(x, y) f_{X, Y}(x, y) \, \mathrm{d} x \, \mathrm{d} y = \iint_{\mathbb{R} ^2} (ax + by) f_{X, Y}(x, y) \, \mathrm{d} x \, \mathrm{d} y \\
    &= a \iint_{\mathbb{R} ^2} x f_{X, Y}(x, y) \, \mathrm{d} x \, \mathrm{d} y + b \iint_{\mathbb{R} ^2} y f_{X, Y}(x, y) \, \mathrm{d} x \, \mathrm{d} y \\
    &= a \int _{\mathbb{R} } x \int_{\mathbb{R} } f_{X, Y}(x, y) \, \mathrm{d} y \, \mathrm{d} x + b \int _{\mathbb{R} } y \int _{\mathbb{R} } f_{X, Y}(x, y) \, \mathrm{d} x \, \mathrm{d} y \\
    &= a \int _{\mathbb{R} } x f_X(x) \, \mathrm{d} x + b \int _{\mathbb{R} } y f_Y(y) \, \mathrm{d} y = a \mathbb{E} [X] + b\mathbb{E} [Y].        
\end{align*}

\paragraph{Application (Monotonicity)} Let \(X, Y\) be random variables on a probability space. If \(X \ge Y\) always, then \(\mathbb{E} [X] \ge \mathbb{E} [Y]\). 

This is because 
\begin{align*}
    X \ge Y \iff X - Y \ge 0 \implies \mathbb{E} [X - Y] \ge 0 \implies \mathbb{E} [X] - \mathbb{E} [Y] \ge 0 \implies \mathbb{E} [X] \ge \mathbb{E} [Y].
\end{align*}

\paragraph{Application (Union Bound/Boole's inequality)}
If \(A_1, A_2, \dots , A_n\) are events, then 
\[
    \mathbb{P} \left( \bigcup_{i=1}^{n} A_i  \right) \le \sum_{i=1}^n \mathbb{P} (A_i).  
\] 
\begin{proof}
    Let \(X_i\) be the indicator for \(A_i\):
    \[
        X_i = \begin{dcases}
            1, &\text{ if } A_i \text{ occurs}  ;\\
            0, &\text{ if } A_i \text{ doesn't occurs}.
        \end{dcases}
    \]  
    Then,
    \[
        \mathbb{E} [X_i] = 1 \cdot \mathbb{P} (A_i) + 0 \cdot \mathbb{P} (A_i^C) = \mathbb{P} (A_i).
    \]
    Let \(X = \sum_{i=1}^n X_i = \#\) of events occur. Let \(Y\) be the indicator for \(\bigcup_{i=1}^{n} A_i \):
    \[
        Y = \begin{dcases}
            1, &\text{ if } X \ge 1 ;\\
            0, &\text{ if } X = 0.
        \end{dcases}
    \]  
    Hence, \(Y \le X\) always. This shows 
    \[
        \mathbb{P} \left( \bigcup_{i=1}^{n} A_i  \right) = \mathbb{E} [Y] \le \mathbb{E} [X] = \mathbb{E} \left[ \sum_{i=1}^n X_i  \right] = \sum_{i=1}^n \mathbb{E} [X_i] = \sum_{i=1}^n \mathbb{P} (A_i).    
    \] 
\end{proof}

\begin{eg}
    We toss a coin \(n\) times. Tosses are independent: Heads with probability \(p\), Tails with probability \(1 - p\). A streak is a maximal subinterval of identical outcomes. What is the expected number of streaks?    
\end{eg}

\begin{explanation}
    Suppose \(X = \#\) of streaks in \(n\) tosses. Let
    \[
        X_i = \begin{dcases}
            1, &\text{ if a new streak starting with } i\text{-th toss};\\
            0, &\text{ if not}.
        \end{dcases}
    \] 
    Then, \(X = \sum_{i=1}^n X_i \) and thus
    \[
        \mathbb{E} [X] = \mathbb{E} \left[ \sum_{i=1}^n X_i  \right] = \sum_{i=1}^n \mathbb{E} [X_i].  
    \] 
    Suppose
    \[
        \mathbb{E} [X_i] = \mathbb{P} (\text{a new streak starting with } i\text{-th toss}) = p_i,
    \]
    then \(p_1 = 1\), and for \(i \ge 2\), 
    \begin{align*}
        p_i &= \mathbb{P} (\text{toss \# } i \neq \text{ toss \# } i - 1) \\
            &= \mathbb{P} \left( \left\{ \text{toss \# } i = H, \text{ toss \#} i - 1 = T   \right\} \cup \left\{ \text{toss \# } i = T, \text{ toss \#} i - 1 = H   \right\}  \right) \\
            &= \mathbb{P} (\text{toss \# } i = H) \mathbb{P} (\text{toss \# } i - 1 = T) + \mathbb{P} (\text{toss \# } i = T) \mathbb{P} (\text{toss \# } i - 1 = H) \\
            &= p(1 - p) + (1 - p)p = 2p(1 - p). 
    \end{align*}  
    Thus, 
    \[
        \mathbb{E} [X] = \sum_{i=1}^n p_i = 1 + 2p(1 - p)(n - 1). 
    \]
\end{explanation}